% T23: 可独立拼入的局部替换；以冻结版行号为准，保留 T21 的第56--70行修正。
% 第88行替换：
上述偶阶电多极相消及单 $j$ 壳层内 $M1\propto\mathbf J$ 的限制，并不排除一切同资数电磁退激。单体奇阶磁多极算符是准自旋标量，在纯单 $j$ 壳层内保持资数，其同资数矩阵元并不因半满而普遍为零。例如，在具备合适能级顺序且单粒子矩阵元非零时，$j=11/2$、$n=6$、$v=2$ 的 $10\to8$ 可有 $M3$ 跃迁：两态同宇称，$|10-8|\le3\le10+8$。跨资数跃迁、组态混合和核芯极化同样可能提供退激贡献；各具体核态的总寿命取决于所有实际开放通道，不由此例确定。

% 第279行替换：
判别式 $\Delta=[(a-b)+(D_1-D_2)]^2+4ab\ge0$，故两根均为实数且不大于零。退化的充要条件为 $ab=0$ 且 $a-b+D_1-D_2=0$，其中包括单向交换的非对角 Jordan 情形；例如 $(a,b,D_1,D_2)=(0,1,2,1)$ 给出 $e^{\mathbf A t}=e^{-2t}\begin{pmatrix}1&t\\0&1\end{pmatrix}$。无源、非负初始布居下，总布居满足 $\dot N_1+\dot N_2=-D_1N_1-D_2N_2\le0$，因此布居有界；此结论来自被动速率结构，不能单凭存在 Jordan 块推断。

% 第293--298行的公式保留，在第298行末加：
此处平衡份额和弛豫时间要求 $a+b>0$；若 $a=b=0$，份额保持初值。
% 第301--302行保留：其一阶快、慢根加权系数正确，渐近展开须 $a+b>0$ 且 D_i/(a+b) 足够小。

% 第309行在 Cramer 法则之前加入：下式仅在 $\det(-\mathbf A)>0$ 时适用。
% 第314行替换：
记 $d=\det(-\mathbf A)=aD_2+bD_1+D_1D_2$。当且仅当 $d>0$ 时，对任意固定非负源 $\mathbf S$ 均存在唯一非负定态，而且从任意初值渐近收敛至该定态；这等价于每个状态均可通过所给有向通道到达某个非零漏失通道。$d=0$ 时，特殊源仍可能与奇异定态方程相容，但解不唯一且不具全局吸引性。例如 $a=b=0$、$D_1=1$、$D_2=0$、$\mathbf S=(1,0)^T$ 有有限定态族 $\mathbf N_{\rm ss}=(1,c)^T$，其中 $c\ge0$。更一般地，令 $M=-\mathbf A$：若 $M\ne0$ 且 $d=0$，给定非负源存在有限非负定态当且仅当 $\operatorname{adj}(M)\mathbf S=0$；若 $M=0$，则当且仅当 $\mathbf S=0$ 时存在定态。特别地，全无漏失且源总和为正时不存在有限定态。

% 第321行最后的 Lindblad 条件句替换为：
非热环境并不自动要求采用 Lindblad 方程；弱耦合 Born--Markov 与久期近似是从具体微观模型导出常见 Markovian Lindblad 方程的充分条件路线，并非 Lindblad 数学形式本身的必要条件。
